% Prozentrechnung – bank/prozentrechnung/ (zusammenbau v0.4, Heft)
\einheitenkopf[t3]{Prozentrechnung}
\setcounter{aufgabe}{9}

\begin{aufgabe}{Umwandeln – Prüfungsaufgaben}
\begin{teile}
\teil $5\,\%$ der Pakete kommen zu spät – welche Aussage passt? Kreuze an. \hfill (P10 2020 OS)\\ \kreuz{$5$ von $10$ Paketen kommen zu spät.}\\ \kreuz{Jedes 5. Paket kommt zu spät.}\\ \kreuz{$5$ von $100$ Paketen kommen zu spät.}
\teil $8\,\%$ der Gäste bestellen Tee – welche Aussage passt? Kreuze an. \hfill (P10 2020 OS)\\ \kreuz{Jeder 8. Gast bestellt Tee.}\\ \kreuz{$8$ von $100$ Gästen bestellen Tee.}\\ \kreuz{$8$ von $10$ Gästen bestellen Tee.}
\teil Jeder vierte Besucher kauft ein Programmheft – wie viel Prozent? Kreuze an. \hfill (P10 2022 OS)\\ \kreuz{$4\,\%$}\\ \kreuz{$25\,\%$}\\ \kreuz{$40\,\%$}\\ \kreuz{$75\,\%$}
\teil Jedes zwanzigste Ei einer Lieferung ist zerbrochen – wie viel Prozent? Kreuze an. \hfill (P10 2022 OS)\\ \kreuz{$2\,\%$}\\ \kreuz{$5\,\%$}\\ \kreuz{$20\,\%$}\\ \kreuz{$80\,\%$}
\end{teile}
\end{aufgabe}

\begin{aufgabe}{Prozentsatz – Prüfungsaufgaben}
\begin{teile}
\teil Eine Stadt hat $40\,000$ Einwohner, $1\,080$ sind in diesem Jahr zugezogen. Wie viel Prozent sind das? \hfill (P10 2014 OS) \\ \leerfeld[\%]
\teil Von $12\,500$ verkauften Konzertkarten wurden $175$ zurückgegeben. Wie viel Prozent sind das? \hfill (P10 2014 OS) \\ \leerfeld[\%]
\end{teile}
\end{aufgabe}

\begin{aufgabe}{Prozentwert – Prüfungsaufgaben}
\begin{teile}
\teil Ein Fernseher kostet $480\,€$. Bei Barzahlung gibt es $15\,\%$ Rabatt. Wie viel Euro spart man? Kreuze an. \hfill (P10 2021 OS)\\ \kreuz{$32\,€$}\\ \kreuz{$408\,€$}\\ \kreuz{$72\,€$}\\ \kreuz{$552\,€$}
\teil Eine Waschmaschine kostet $650\,€$. Im Angebot gibt es $30\,\%$ Rabatt. Wie viel Euro spart man? Kreuze an. \hfill (P10 2021 OS)\\ \kreuz{$455\,€$}\\ \kreuz{$65\,€$}\\ \kreuz{$845\,€$}\\ \kreuz{$195\,€$}
\teil Ein Tablet kostet $240{,}00\,€$. Es gibt $15\,\%$ Rabatt. Wie viel Euro Rabatt sind das? \hfill (P10 2017 OS) \\ \leerfeld[€]
\teil Eine Gitarre kostet $180{,}00\,€$. Es gibt $35\,\%$ Rabatt. Wie viel Euro Rabatt sind das? \hfill (P10 2017 OS) \\ \leerfeld[€]
\teil $17\,\%$ von $60\,€$ \hfill (P10 2014 OS) \\ \leerfeld[€]
\teil $19\,\%$ von $40\,€$ \hfill (P10 2014 OS) \\ \leerfeld[€]
\teil $40\,\%$ von $85\,€$ \hfill (P10 2026 FOR) \\ \leerfeld[€]
\teil $60\,\%$ von $45\,€$ \hfill (P10 2026 FOR) \\ \leerfeld[€]
\end{teile}
\end{aufgabe}

\begin{aufgabe}{Grundwert – Prüfungsaufgaben}
\begin{teile}
\teil Beim Kauf einer Mütze spart Mira $4\,€$, das sind $25\,\%$ des alten Preises. Wie viel kostete die Mütze vorher? \hfill (P10 2025 OS) \\ \leerfeld[€]
\teil Bei einer Konzertkarte beträgt der Rabatt $15\,€$, das sind $30\,\%$ des alten Preises. Wie hoch war der alte Preis? \hfill (P10 2025 OS) \\ \leerfeld[€]
\teil Lisa bekommt $40\,\%$ einer Erbschaft, das sind $6\,000\,€$. Wie groß ist die ganze Erbschaft? \hfill (P10 2023 OS) \\ \leerfeld[€]
\teil $20\,\%$ der Ernte eines Obstbauern sind $450$ kg Äpfel. Wie groß ist die ganze Ernte? \hfill (P10 2023 OS) \\ \leerfeld[kg]
\end{teile}
\end{aufgabe}

\begin{aufgabe}{Veränderung – Prüfungsaufgaben}
\begin{teile}
\teil Ein Busticket kostet $2{,}40\,€$ und wird um $25\,\%$ teurer. Wie viel kostet es dann? \hfill (P10 2024 OS) \\ \leerfeld[€]
\teil Ein Schwimmkurs kostet $65\,€$ und wird um $12\,\%$ teurer. Wie viel kostet er dann? \hfill (P10 2024 OS) \\ \leerfeld[€]
\end{teile}
\end{aufgabe}
